A classifier deployed in practice sees inputs that differ from its training data, and the usefulness of its probabilities then depends on whether they remain honest. Post-hoc temperature scaling is the cheapest calibration fix~\cite{guo2017calibration}, but it is fitted on data that look like the training distribution. Whether the fitted temperature transfers to shifted inputs is an empirical question that large-scale benchmarks have studied on image corruptions~\cite{ovadia2019can} (building on the corruption benchmark of Hendrycks and Dietterich~\cite{hendrycks2019benchmarking}), with the conclusion that in-domain calibration does not carry over reliably~\cite{ovadia2019can,tomani2020post}.

We ask the same question at a scale where every cell of the experiment can be run in seconds on a CPU and repeated over five seeds: scikit-learn digits (8$\times$8 images), an MLP, and two shift families with a graded intensity, rotation and additive Gaussian noise. We compare a single MLP, TS fitted on in-distribution validation data, MC dropout~\cite{gal2015dropout} and a five-member deep ensemble~\cite{lakshminarayanan2016simple}, measured by accuracy, expected calibration error (ECE), negative log-likelihood (NLL) and Brier score. The question that motivated the study is: \emph{does temperature scaling fitted in-distribution stay calibrated under shift?} We register five hypotheses in advance (Section~\ref{sec:setup}), and add an oracle-temperature ablation and two sensitivity sweeps to probe the mechanism.

Contributions are modest by design: (i) a small-scale reproduction of the shift-calibration comparison on a different data modality and two shift families; (ii) a controlled ablation that separates ``the temperature is wrong under shift'' from ``a scalar cannot fix shifted predictions''; (iii) all code, run logs and tables in the repository. Everything is reimplemented here, nothing is tuned, and the conclusions are limited to this dataset, architecture and shift construction.
